\subsection{Bounded input batches and two-pass construction}

Rotating all $N$ raw vectors at once requires two FP32 device arrays,
together $8Nd$ bytes. The builder instead stages bounded batches through
pinned host memory and device raw/rotated buffers. For $n_b$ pipeline
buffers, configured add rows $a_{\rm add}$ and per-array byte budget
$S_{\rm stage}$, the source stages
\begin{equation}
 a_{\rm in}=\max\!\left\{1,\min\!\left(a_{\rm add},
 \left\lfloor\frac{S_{\rm stage}}{4d}\right\rfloor\right)\right\},
 \qquad S_{\rm float}=8n_ba_{\rm in}d.
 \label{eq:4-staging-rows}
\end{equation}
Three buffers and the default 128\,MiB budget give at most 768\,MiB of
input/rotated device floats. Assignment products and encoder scratch
are separate allocations.

One-pass add encodes input-order arrays, then gathers them into list order, temporarily keeping both copies. Two-pass add assigns/sorts IDs first, then revisits input to encode and scatter directly. It trades another read/rotation for lower simultaneous storage. Raw-input assignment against back-transformed centroids uses the real-arithmetic orthogonality identity.

\subsection{Phase occupancy and the process peak}

The one-pass capacity estimate is
\begin{equation}
 H_{\rm single}=N(2b_r+3M+20)+2^{30}\quad\text{bytes}.
 \label{eq:4-peak-heuristic}
\end{equation}
Ten live gather/transpose arrays contribute $3NM+2Nb_r+20N$ bytes: three primary copies, two residual/correction copies, assignments, sort order and IDs. The extra 1\,GiB is heuristic headroom, not a full peak model.

At final two-pass repack, IDs, residuals, corrections and both primary layouts occupy
\begin{equation}
 S_{\rm repack,2}=N(b_r+2M+8).
 \label{eq:4-two-pass-repack}
\end{equation}
This counts that phase only. Whole-build peak also considers earlier pipeline/inverse scratch, model, sort, context and later workspace allocations.

\begin{table}[tbp]
\centering
\caption[Analytical row-array occupancy by construction phase]{Analytical row-array occupancy (GiB), $d=1{,}024$, $M=128$,
$B_r=8$. BGE uses one pass; Stella uses two. The one-pass Stella column
is counterfactual capacity analysis. Pipeline, model, offsets and runtime
costs are excluded from these phase values.}
\label{tab:4-phase-storage}
\small
\begin{tabular}{@{}lrrr@{}}
\toprule
Dataset & Final row payload & One-pass gather & Two-pass repack \\
\midrule
BGE & 10.90 & 23.05 & 12.11 \\
Stella & 19.20 & 40.59 & 21.32 \\
\bottomrule
\end{tabular}
\end{table}

\Cref{tab:4-phase-storage} quantifies the compressed-copy saving. Stella's
cold log selects two passes and prints 41.6\,GiB for
\cref{eq:4-peak-heuristic}, against 30.6\,GiB available at the tail.
The separately sampled resident process totals, including context and
query workspace, are 11.93\,GiB for BGE and 20.44\,GiB for Stella
(\cref{sec:4-4}). The supplementary 20\,ms device samples observe a 26.08\,GiB BGE peak
near the start of add, and a 22.58\,GiB Stella peak during repack.
These are sampled device totals including context; shorter-lived peaks
can be missed.

Assignment-product capacity uses an independent configured row count
$a_{\rm assign}$, selected by \texttt{JHQ\_ASSIGN\_BATCH} or the source
default: 32,768 for at least 8,192 lists, otherwise 8,192.
\begin{equation}
 S_{\rm dots}=4n_b n_{\rm list}a_{\rm assign}\quad\text{bytes}.
 \label{eq:4-assignment-scratch}
\end{equation}
Allocation capacity need not shrink with $a_{\rm in}$ or the rows in a
particular call. \Cref{tab:b-scratch} checks four defaults: the 1,536/3,072-dimensional
cases allocate 3\,GiB, despite different staging rows. BGE/Stella still
allocate three 4\,GiB arrays, or 12\,GiB, near add entry. The 23.05\,GiB gather
payload describes a later lifetime: adding 12\,GiB to it would not give
the actual peak. Peak modelling must take the maximum over stage-specific
live arrays, including pipeline scratch. Stella's smaller two-pass
repack payload, plus model/context, explains its lower observed peak
despite its larger database.

\subsection{Construction measurements}

